\documentclass[11pt]{article}
\usepackage{preamble}
\setlength{\parskip}{4pt}

\begin{document}

\daybanner{AP Statistics \textperiodcentered\ Topic 6.8 (CED 3.10) \textperiodcentered\ B: 2026-12-11 \textperiodcentered\ E: 2027-01-22 \textperiodcentered\ TEACHER KEY}{One Difference, Not Two Separate Intervals}

\sectionbanner{CED TETHER}

{\small
\begin{itemize}[leftmargin=1.5em,itemsep=2pt,topsep=2pt]
  \item \textbf{Skill 1.D} --- Identify an appropriate inference method for confidence intervals.
  \item \textbf{Skill 4.C} --- Verify that inference procedures apply in a given situation.
  \item \textbf{Skill 3.D} --- Construct a confidence interval, provided conditions for inference are met.
  \item \textbf{EU UNC-4} --- An interval of values should be used to estimate parameters, in order to account for uncertainty.
  \item \textbf{LO UNC-4.I} --- Identify an appropriate confidence interval procedure for a comparison of population proportions. [Skill 1.D]
  \item \textbf{EK UNC-4.I.1} --- The appropriate confidence interval procedure for a two-sample comparison of proportions for one categorical variable is a two-sample z-interval for a difference between population proportions.
  \item \textbf{LO UNC-4.J} --- Verify the conditions for calculating confidence intervals for a difference between population proportions. [Skill 4.C]
  \item \textbf{EK UNC-4.J.1} --- In order to calculate confidence intervals to estimate a difference between proportions, we must check for independence and that the sampling distribution is approximately normal: a. To check for independence: i. Data should be collected using two independent, random samples or a randomized experiment.
  \item ii. When sampling without replacement, check that n1 $\leq$ 10\%N1 and n2 $\leq$ 10\%N2.
  \item b. To check that the sampling distribution of $\widehat p_1$ - $\widehat p_2$ is approximately normal (shape): i. For categorical variables, check that n1*$\widehat p_1$, n1*(1-$\widehat p_1$), n2*$\widehat p_2$, and n2*(1-$\widehat p_2$) are all greater than or equal to some predetermined value, typically either 5 or 10.
  \item \textbf{LO UNC-4.K} --- Calculate an appropriate confidence interval for a comparison of population proportions. [Skill 3.D]
  \item \textbf{EK UNC-4.K.1} --- For a comparison of proportions, the interval estimate is ($\widehat p_1$ - $\widehat p_2$) $\pm$ z* $\cdot$ sqrt($\widehat p_1$(1-$\widehat p_1$)/n1 + $\widehat p_2$(1-$\widehat p_2$)/n2). (Note: these formulas can be constructed from the general formula and SE formulas on the AP formula sheet; they need not be memorized.)
  \item \textbf{LO UNC-4.L} --- Calculate an interval estimate based on a confidence interval for a difference of proportions. [Skill 3.D]
  \item \textbf{EK UNC-4.L.1} --- Confidence intervals for a difference in proportions can be used to calculate interval estimates with specified units.
\end{itemize}
}
\textbf{Essential question:} How do the study design and variance structure determine a valid interval for a difference in population proportions?\par
\textbf{DOK-3 skill:} 4.C \textperiodcentered\ \textbf{FRQ pattern:} {\small\texttt{two-\allowbreak{}sample-\allowbreak{}interval-\allowbreak{}vs-\allowbreak{}separate-\allowbreak{}intervals}}\par
\textbf{DOK rationale:} Part (c) requires adjudicating two plausible interval methods by coordinating the two-sample design, condition evidence, variance structure, confidence-level target, and the inferential consequence of an invalid construction.\par

\frameworkphaseheader{Do Now (first take) $\rightarrow$ Explore (video + follow-along) $\rightarrow$ Exit (finish + turn in)}{1 $\rightarrow$ 3}{45}{%
\begin{itemize}[leftmargin=1.5em,itemsep=2pt,topsep=2pt]
  \item Board slide up at the bell. Ask students to choose a method before confirming it.
  \item Begin the video at minute 5; return at minute 33 to separate variance, SE, and margin of error.
  \item Collect at the door. Score part (c) only, E/P/I.
\end{itemize}
}{%
\begin{itemize}[leftmargin=1.5em,itemsep=2pt,topsep=2pt]
  \item Choose Elena or Devon and cite design, counts, or zero.
  \item Complete the follow-along, finish (a)--(c), revise, and turn in the sheet.
\end{itemize}
}{%
\begin{itemize}[leftmargin=1.5em,itemsep=2pt,topsep=2pt]
  \item What single parameter is the town comparison trying to estimate?
  \item Which four observed counts justify the normal approximation?
  \item Do standard errors or variances add?
\end{itemize}
}{Solo teacher. Circulate ~5 pairs. Link the formula to design and the rejected method to zero.}

\begin{callout}{\IconWarn\ WATCH FOR}{calloutred}
\begin{itemize}[leftmargin=1.5em,itemsep=2pt,topsep=2pt]
  \item Pooling the two sample proportions even though the task is to estimate an unrestricted difference.
  \item Adding separate margins of error instead of adding estimated variance terms under the square root.
  \item Assuming two individual 95 percent intervals automatically yield a 95 percent interval for their difference.
\end{itemize}
\end{callout}

\sectionbanner{SCORING PART (c) --- E / P / I}

\textbf{Expected elements}
\begin{itemize}[leftmargin=1.5em,itemsep=2pt,topsep=2pt]
  \item \textbf{selects-two-sample-method} (required) --- Selects Elena's two-sample $z$-interval
  \item \textbf{uses-design-and-counts} (required) --- Uses independent SRSs and counts 77, 63, 48, and 72
  \item \textbf{explains-variance-combination} (required) --- Adds independent variance terms under one square root
  \item \textbf{rejects-separate-interval-coverage} (required) --- Rejects automatic 95 percent coverage from two intervals
  \item \textbf{traces-inferential-consequence} (required) --- Explains Devon's misleading inclusion of zero
\end{itemize}

{\small\renewcommand{\arraystretch}{1.25}
\noindent\begin{tabular}{|p{0.5in}|p{5.9in}|}
\hline
\textbf{E} & Selects Elena, uses design and counts, explains variance and coverage, and diagnoses the false zero conclusion. \\ \hline
\textbf{P} & Selects Elena but incompletely explains the rejected construction or its consequence. \\ \hline
\textbf{I} & Uses separate intervals, pools samples, or ignores design and zero. \\ \hline
\end{tabular}}

\clearpage
\focusbanner{TODAY'S PROBLEM --- annotated}

Town Oak has 2{,}200 households and Pine 1{,}800. Independent SRSs find 77 of 140 Oak households and 48 of 120 Pine households use curbside composting. Let $p_O-p_P$ mean Oak minus Pine.\par\smallskip \textbf{Elena:} ``The two-sample $z$-interval is $(0.030,0.270)$.''\par \textbf{Devon:} ``Subtract the farthest endpoints of separate 95\% intervals. That gives $(-0.020,0.320)$, so zero is plausible.''\par
\begin{center}
\textbf{Household samples}\par\smallskip
\begin{tabular}{|l|c|c|c|c|}
\hline
\textbf{Town} & \textbf{Use} & \textbf{Other} & \textbf{n} & \textbf{N} \\ \hline
\textbf{Oak} & 77 & 63 & 140 & 2200 \\ \hline
\textbf{Pine} & 48 & 72 & 120 & 1800 \\ \hline
\end{tabular}
\end{center}
\begin{firsttakebox}{FIRST TAKE (5 min)}
Which method gives the requested 95\% interval for $p_O-p_P$? Cite one design or numerical detail.\par
\answer{Not graded. Equal centers do not make both uncertainty calculations valid.}
\end{firsttakebox}

\textit{Rules callout ``BUILD ONE INTERVAL FOR THE DIFFERENCE (keep on the page)'' is printed on the student sheet.}\par\medskip

\dokbadge{1}~\textbf{(a)} Define $p_O$ and $p_P$ and name the interval procedure for $p_O-p_P$.\par
\answer{$p_O$ and $p_P$ are the proportions of all Oak and Pine households using their town's service. Use a two-sample $z$-interval for $p_O-p_P$.}

\dokbadge{2}~\textbf{(b)} Check independent samples, both 10 percent conditions, and all four counts. Compute the estimate, SE, margin of error, and 95\% interval.\par
\answer{The SRSs are independent; $140\le220$ and $120\le180$; and counts 77, 63, 48, and 72 exceed 10. Thus $\hat p_O-\hat p_P=0.55-0.40=0.15$, $SE=0.0613829$, and $ME=1.96(0.0613829)=0.1203104$. The interval is $(0.0296896,0.2703104)\approx(0.030,0.270)$.}

\dokbadge{3}~\textbf{(c)} Choose the warranted method. Use design, counts, variance structure, and zero's location. Explain why separate 95\% intervals do not produce the requested coverage and how Devon's result misleads.\par
\answer{Elena is warranted. Independent SRSs and counts 77, 63, 48, and 72 support one two-sample interval. Each sample proportion estimates its own population proportion; independence makes the two estimated variance terms add under one square root. Devon instead combines separate intervals $(0.4676,0.6324)$ and $(0.3123,0.4877)$ into $(-0.0201,0.3201)$. Two individual 95\% intervals do not guarantee one 95\% interval for their difference. His wider result includes zero and falsely makes no difference look plausible; Elena's requested interval supports a positive Oak-minus-Pine difference of about 3 to 27 points.}

\begin{summaryexitbox}{EXIT}
One thing the video changed about my first take: \hrulefill
\end{summaryexitbox}

\end{document}
